\chapter{Contexte et données}
\label{chap:contexte}

\section{JeuxDeMots}

\jdm{} est un jeu en ligne sérieux développé par le LIRMM (Montpellier) dont le but est de construire collaborativement un réseau lexical du français. Les joueurs associent des mots à des propriétés, des catégories, des synonymes, des actions possibles, etc. Le résultat est un graphe de connaissances contenant des millions de nœuds et de relations.

\section{Structure du graphe JDM}

Chaque \textbf{nœud} représente un concept (mot, expression, syntagme) et est identifié par un entier unique (ID JDM). Chaque \textbf{relation} est un triplet (nœud source, type de relation, nœud cible) associé à un poids (score de confiance) allant de 1 à plusieurs milliers.

Les types de relations principaux utilisés dans \wido{} sont :

\begin{table}[H]
    \centering
    \begin{tabular}{lll}
        \toprule
        Nom & ID & Signification \\
        \midrule
        \rel{r\_isa} & 6 & Appartient à la catégorie \\
        \rel{r\_has\_part} & 9 & Possède la partie \\
        \rel{r\_can\_eat} & 43 & Peut manger \\
        \rel{r\_color} & 42 & A pour couleur \\
        \rel{r\_carac} & 12 & A pour caractéristique \\
        \rel{r\_synt\_synon} & 1 & Synonyme syntaxique \\
        \bottomrule
    \end{tabular}
    \caption{Types de relations JDM utilisés dans WIDO}
    \label{tab:relations}
\end{table}

\section{API JDM v0}

L'API REST publique du LIRMM expose les données JDM via des endpoints structurés. Elle supporte la pagination native (\code{limit} + \code{offset}) et le filtrage par type de relation (\code{types\_ids}) et poids minimum (\code{min\_weight}). Le tableau~\ref{tab:endpoints} résume les endpoints exploités dans \wido{}.

\begin{table}[H]
    \centering
    \begin{tabular}{ll}
        \toprule
        Endpoint & Usage \\
        \midrule
        \code{/v0/node\_by\_name/\{name\}} & Résoudre un nom en nœud JDM \\
        \code{/v0/relations/from/\{name\}} & Récupérer les relations sortantes \\
        \code{/v0/relations/to/\{name\}} & Récupérer les relations entrantes \\
        \code{/v0/relations/from\_by\_id/\{id\}} & Relations sortantes par ID \\
        \code{/v0/relations/to\_by\_id/\{id\}} & Relations entrantes par ID \\
        \code{/v0/relations/from\_by\_id/\{id1\}/to\_by\_id/\{id2\}} & Vérification d'un couple \\
        \bottomrule
    \end{tabular}
    \caption{Endpoints API JDM utilisés par WIDO}
    \label{tab:endpoints}
\end{table}
